%%%%%%%%%%%%%%%%%%%%%%%%%%%%%%%%%%%%%%%%%%%%%%%%%%%%%%%%%%%%%%%%%%%%%%%%%%%%%%%
% NSF TTP-T Proposal Template for NxTAGE/Aria
% Based on NSF 25-540 Solicitation Requirements
% 
% Key Constraints:
%   - Project Description: 15 pages maximum
%   - Executive Summary: 1 page maximum  
%   - Budget: Up to $1,200,000 for up to 36 months
%   - I-Corps training required ($50,000 in Year 1)
%%%%%%%%%%%%%%%%%%%%%%%%%%%%%%%%%%%%%%%%%%%%%%%%%%%%%%%%%%%%%%%%%%%%%%%%%%%%%%%

\documentclass[11pt]{article}

% Page setup per NSF requirements
\usepackage[margin=1in]{geometry}
\usepackage{times}  % NSF requires Times New Roman or similar
\usepackage{graphicx}
\usepackage{booktabs}
\usepackage{hyperref}
\usepackage{enumitem}
\usepackage{xcolor}
\usepackage{fancyhdr}
\usepackage{titlesec}

% Section formatting
\titleformat{\section}{\normalfont\Large\bfseries}{\thesection.}{1em}{}
\titleformat{\subsection}{\normalfont\large\bfseries}{\thesubsection}{1em}{}
\titleformat{\subsubsection}{\normalfont\normalsize\bfseries}{\thesubsubsection}{1em}{}

% Comments for drafting (remove for final)
\newcommand{\TODO}[1]{\textcolor{red}{[TODO: #1]}}
\newcommand{\NOTE}[1]{\textcolor{blue}{[NOTE: #1]}}
\newcommand{\CITE}[1]{\textcolor{orange}{[CITE: #1]}}

% Header/Footer
\pagestyle{fancy}
\fancyhf{}
\rhead{NSF TTP-T: Aria}
\lhead{Project Description}
\rfoot{Page \thepage}

\begin{document}

%%%%%%%%%%%%%%%%%%%%%%%%%%%%%%%%%%%%%%%%%%%%%%%%%%%%%%%%%%%%%%%%%%%%%%%%%%%%%%%
% COVER PAGE (for internal use - not part of 15-page limit)
%%%%%%%%%%%%%%%%%%%%%%%%%%%%%%%%%%%%%%%%%%%%%%%%%%%%%%%%%%%%%%%%%%%%%%%%%%%%%%%

% \begin{center}
% \LARGE\textbf{NSF TTP-T: Aria -- Privacy-Preserving Personal AI \\
% Through Dynamic Persona Modeling}

% \vspace{1em}
% \large
% Principal Investigator: [PI Name]\\
% Institution: University of Washington\\
% Requested Amount: \$[Amount]\\
% Duration: 36 months

% \vspace{2em}
% \textit{Internal Draft -- Not for Submission}
% \end{center}

% \newpage

%%%%%%%%%%%%%%%%%%%%%%%%%%%%%%%%%%%%%%%%%%%%%%%%%%%%%%%%%%%%%%%%%%%%%%%%%%%%%%%
% PROJECT SUMMARY (Separate document, 1 page)
% NOTE: This is a separate submission component, not part of 15-page limit
%%%%%%%%%%%%%%%%%%%%%%%%%%%%%%%%%%%%%%%%%%%%%%%%%%%%%%%%%%%%%%%%%%%%%%%%%%%%%%%

%\pagenumbering{gobble}
\thispagestyle{empty}
\section*{PROJECT SUMMARY}

\subsection*{Overview}
% If PI has recent I-Corps on this technology, list award number on first line
% Example: NSF I-Corps Award \#XXXXXXX (if applicable for waiver)

\TODO{Write 1 paragraph covering:}
\begin{itemize}[noitemsep]
    \item Societal need for the proposed technology translation
    \item Areas of application (personal AI assistants, productivity tools)
    \item Potential stakeholders (consumers, enterprises, device manufacturers)
    \item Potential outcomes (startup formation, licensing, product deployment)
\end{itemize}

\textit{[Draft: Current personal AI assistants fail to provide truly personalized 
assistance because they either sacrifice user privacy through cloud-based processing 
or lack the contextual understanding needed for meaningful personalization. Aria 
addresses this gap through breakthrough ``Sims Technology'' that creates dynamic, 
evolving digital representations of user preferences while keeping all processing 
on-device. This project will translate foundational research in privacy-preserving 
machine learning and personalization into a commercial product targeting the \$XX 
billion personal assistant market...]}

\subsection*{Intellectual Merit}
\TODO{Write 1 paragraph addressing:}
\begin{itemize}[noitemsep]
    \item Goals of the proposed research
    \item Technical hurdles to be addressed
    \item Novel contributions to the field
\end{itemize}

\textit{[Draft: The project advances knowledge in on-device machine learning, 
privacy-preserving personalization, and multi-context user modeling. Key technical 
challenges include developing efficient fine-tuning algorithms that operate within 
mobile device constraints, creating robust persona synthesis methods that blend 
multiple user contexts, and building secure architectures that prevent data 
leakage while enabling meaningful AI assistance...]}

\subsection*{Broader Impacts}
\TODO{Write 1 paragraph describing:}
\begin{itemize}[noitemsep]
    \item Expected commercial outcomes
    \item Economic benefits
    \item Societal impacts (privacy, accessibility, etc.)
\end{itemize}

\textit{[Draft: Successful translation will establish a new paradigm for 
privacy-respecting AI assistance, potentially disrupting the current 
cloud-dependent model. Commercial outcomes include startup formation and 
strategic partnerships with device manufacturers. The technology promotes 
user data sovereignty and addresses growing concerns about AI privacy...]}

\newpage

%%%%%%%%%%%%%%%%%%%%%%%%%%%%%%%%%%%%%%%%%%%%%%%%%%%%%%%%%%%%%%%%%%%%%%%%%%%%%%%
% PROJECT DESCRIPTION (15 pages maximum)
%%%%%%%%%%%%%%%%%%%%%%%%%%%%%%%%%%%%%%%%%%%%%%%%%%%%%%%%%%%%%%%%%%%%%%%%%%%%%%%

\setcounter{page}{1}

%\begin{center}
%{\bf \Large {NSF TTP-T: Aria -- Privacy-Preserving Personal AI \\
% Through Dynamic Persona Modeling}\\
%\large{Project Description}
%}
%\end{center}

%%%%%%%%%%%%%%%%%%%%%%%%%%%%%%%%%%%%%%%%%%%%%%%%%%%%%%%%%%%%%%%%%%%%%%%%%%%%%%%
% Section 1: Executive Summary
% NSF TTP-T Proposal: NxTAGE/Aria Privacy-First Personal AI Assistant
% Target length: ~1 page
%%%%%%%%%%%%%%%%%%%%%%%%%%%%%%%%%%%%%%%%%%%%%%%%%%%%%%%%%%%%%%%%%%%%%%%%%%%%%%%

\section{Executive Summary}
\label{sec:summary}
\vspace*{-12pt}
\subsection{Identification of Drivers}\vspace*{-8pt}

Contemporary AI-powered personal assistants require users to surrender sensitive data to cloud infrastructure in exchange for personalization---a trade-off that research shows creates significant adoption barriers. Studies indicate that approximately 60\% of voice assistant users express privacy concerns, while nearly half remain unaware that devices continuously monitor ambient audio \cite{edu2021security, li2023security}. This privacy-personalization paradox excludes entire market segments: privacy-conscious consumers avoid existing solutions, and organizations in regulated industries (healthcare, legal, financial services) cannot deploy cloud-dependent assistants due to compliance requirements under HIPAA, GDPR, and similar frameworks.

Translation is needed because no commercial solution currently delivers meaningful personalization while guaranteeing complete data sovereignty. Existing privacy-preserving approaches---including federated learning and limited on-device processing---still involve server-side coordination or sacrifice personalization depth. The intelligent virtual assistant market, projected to reach \$14--80 billion by 2030 at 24\%+ CAGR \cite{grandview2023iva}, presents substantial opportunity for solutions that resolve this fundamental tension rather than merely mitigating it.

\vspace*{-10pt}
\subsection{The Innovation}
\vspace*{-8pt}
Aria introduces \textbf{Sims Technology}---dynamic digital representations that model user preferences across distinct life contexts while processing all data exclusively on-device. Unlike existing assistants that maintain monolithic user profiles, Sims Technology recognizes that users behave differently across professional, personal, and social contexts, automatically adapting interaction patterns without explicit switching or invasive monitoring.
Three technical innovations enable this capability. First, \textbf{on-device personalization} leverages Low-Rank Adaptation (LoRA), which reduces trainable parameters by approximately 10,000-fold compared to full fine-tuning \cite{hu2022lora}, enabling sophisticated model adaptation within mobile device constraints. Second, \textbf{multi-context persona synthesis} employs sensor fusion techniques achieving greater than 90\% accuracy in context detection \cite{wang2019deep} to automatically select appropriate behavioral models. Third, \textbf{privacy-preserving architecture} implements differential privacy guarantees using techniques proven effective for deep neural network training \cite{abadi2016deep}, with hardware-level protection through ARM TrustZone Trusted Execution Environments \cite{pinto2019demystifying}.

This approach is transformative because it eliminates rather than manages the privacy-personalization trade-off. Where competitors offer privacy controls or data minimization, Aria ensures sensitive information never leaves the device under any circumstance---a fundamental architectural distinction that enables deployment in contexts where existing solutions cannot operate.

\vspace*{-10pt}
\subsection{The Value Proposition}
\vspace*{-8pt}
Aria delivers differentiated value across three stakeholder groups. \textit{Privacy-conscious consumers} gain AI assistance matching cloud-based alternatives without data exposure, addressing documented concerns that currently suppress adoption. \textit{Enterprises in regulated industries} can deploy AI productivity tools previously prohibited by compliance requirements, unlocking substantial efficiency gains. \textit{Device manufacturers} obtain licensing opportunities for privacy-differentiated AI capabilities as hardware commoditization intensifies competitive pressure.
The key differentiator is \textbf{complete on-device processing with multi-context awareness}---no competitor combines these capabilities. We measure differentiation through: (1) technical benchmarks comparing personalization quality against cloud-based alternatives under equivalent computational budgets; (2) user trust metrics validated through longitudinal studies; and (3) market validation demonstrating enterprise adoption in previously inaccessible regulated sectors.

Societal value extends beyond commercial outcomes. By demonstrating viable alternatives to surveillance-based AI business models, Aria contributes to rebuilding public trust in AI systems---surveys consistently rank privacy among the most important factors for AI trust \cite{gillespie2023trust}. The technology enables beneficial AI applications in sensitive domains currently underserved due to privacy barriers.

%A provisional patent application filed in 2025 protects core Sims Technology innovations. This project establishes the translation pathway from validated research to commercial deployment through I-Corps customer discovery, iterative prototype development, and market validation with target user populations.

% Optional for TTP-T, include if you have partners
% \subsection{NSF-Catalyzed Partnership}
% \TODO{If applicable, describe partner collaboration}

%%%%%%%%%%%%%%%%%%%%%%%%%%%%%%%%%%%%%%%%%%%%%%%%%%%%%%%%%%%%%%%%%%%%%%%%%%%%%%%
% Section 2: Research Motivation
% NSF TTP-T Proposal for NxTAGE/Aria
% Suggested length: 3-5 pages (this version targets 3-4 pages)
%%%%%%%%%%%%%%%%%%%%%%%%%%%%%%%%%%%%%%%%%%%%%%%%%%%%%%%%%%%%%%%%%%%%%%%%%%%%%%%

\vspace*{-8pt}
\section{Research Motivation}
\label{sec:motivation}

\vspace*{-8pt}
\subsection{Problem Statement and Importance}
\vspace*{-8pt}
\textbf{The Privacy-Personalization Paradox.} Contemporary AI-powered personal assistants confront a fundamental architectural tension: delivering meaningful personalization typically requires transmitting sensitive user data to centralized cloud infrastructure for processing and model training. This design creates an inherent conflict between user privacy expectations and the data requirements for effective AI assistance. Research demonstrates that users exhibit significant apprehension regarding voice assistant data practices, with studies indicating that approximately 60\% of users express at least occasional concerns about privacy when interacting with these systems \cite{li2023security}. Furthermore, investigations reveal that nearly half of users remain unaware that voice assistants continuously monitor ambient audio for activation commands, exacerbating trust deficits when this functionality becomes apparent \cite{edu2021security}.

This tension manifests across several critical dimensions. First, regarding \textbf{user trust and adoption barriers}, privacy concerns represent a primary factor inhibiting technology acceptance. Research on voice-based digital assistants indicates that perceived privacy risk significantly influences adoption intentions, with users engaging in privacy calculus behaviors that weigh personalization benefits against data exposure costs \cite{vimalkumar2021digital}. Second, concerning \textbf{regulatory and compliance pressure}, the global regulatory landscape increasingly constrains centralized data collection practices. The European General Data Protection Regulation (GDPR), California Consumer Privacy Act (CCPA), and emerging international frameworks impose stringent requirements on data processing, storage, and user consent that cloud-dependent architectures struggle to satisfy efficiently. Third, with respect to \textbf{security vulnerabilities}, centralized repositories of personal data present attractive targets for malicious actors. Documented incidents include unintended recording transmission to contacts, employee access to private conversations for algorithm training, and various attack vectors exploiting voice assistant vulnerabilities \cite{edu2021security}. Finally, regarding \textbf{context collapse in single-persona modeling}, existing cloud-based systems represent users through monolithic preference profiles that fail to capture the multifaceted nature of human identity. Users behave differently across life contexts---professional, familial, recreational---yet current assistants cannot distinguish or adapt to these contextual variations without invasive monitoring that further compounds privacy concerns.

\textbf{Stakeholder Impact and Market Need.} The consequences of this privacy-personalization paradox extend across the AI assistant ecosystem. \textit{Consumers} face an unsatisfactory choice between accepting generic, context-unaware assistance or surrendering personal data sovereignty. \textit{Enterprise organizations} in regulated industries (healthcare, finance, legal) cannot deploy powerful AI assistance tools without violating compliance requirements governing sensitive data handling. \textit{Device manufacturers} seek differentiation strategies as hardware commoditization intensifies, yet remain dependent on third-party cloud services for AI capabilities. \textit{Healthcare and financial services providers} recognize the transformative potential of AI assistance but cannot reconcile data protection obligations with current architectural approaches.

\vspace*{-8pt}
\subsection{Current State-of-the-Art and Gaps}
\vspace*{-8pt}
The personal AI assistant landscape comprises several architectural approaches, each exhibiting significant limitations that Aria's innovation addresses.

\textbf{Cloud-Centric Assistants.} Dominant platforms including Apple Siri, Amazon Alexa, and Google Assistant process user interactions primarily through cloud infrastructure. While this approach enables access to substantial computational resources, it fundamentally compromises user privacy by transmitting voice recordings, queries, and contextual data to remote servers. Research characterizing the security and privacy challenges of virtual assistants identifies persistent concerns regarding recording practices, data retention policies, and the potential for both intentional misuse and unintended exposure \cite{edu2021security}. These systems typically employ single-user-profile architectures that aggregate preferences without distinguishing contextual variations, limiting personalization depth and appropriateness.

\textbf{Privacy-Preserving Approaches.} Emerging alternatives attempt to address privacy concerns through various mechanisms. Federated learning, pioneered for mobile keyboard prediction, enables collaborative model training without centralizing raw data by exchanging model updates rather than training examples \cite{mcmahan2017communication}. This paradigm shift demonstrated that effective learning could occur with data remaining distributed across devices, reducing privacy exposure. However, federated approaches still involve server-side aggregation and do not address the fundamental challenge of personalization requiring intimate knowledge of individual user contexts. On-device processing initiatives, including Apple's expansion of Neural Engine capabilities, represent movement toward local computation but remain limited in scope and personalization depth.

\textbf{Browser-Dependent Solutions.} Recent entrants such as browser-based AI assistants leverage local storage and processing within browser environments. While partially addressing privacy concerns for web-based activities, these solutions cannot operate across the full spectrum of user contexts---they remain blind to offline activities, non-browser applications, and the physical environment signals that inform genuine context awareness.

\textbf{Critical Gaps.} Analysis of current approaches reveals several unaddressed requirements. No existing solution provides \textit{multi-persona modeling} that dynamically represents and transitions between user contexts. Current privacy-preserving techniques sacrifice \textit{personalization depth} for data protection rather than achieving both simultaneously. Existing systems lack \textit{evolutionary adaptation} mechanisms that allow user representations to develop organically over time without explicit training interventions. Finally, no commercial solution offers \textit{complete data sovereignty} where user information never leaves device boundaries under any circumstance.

\vspace*{-8pt}
\subsection{The STEM Discovery Foundation}
\vspace*{-8pt}
Aria's innovation synthesizes advances across multiple research domains into a coherent architectural approach that resolves the privacy-personalization paradox.

\textbf{On-Device Machine Learning.} The past decade witnessed substantial progress in deploying machine learning models on resource-constrained edge devices. Research on deploying AI at the network periphery demonstrates that optimized inference frameworks, model compression techniques, and specialized hardware accelerators now enable sophisticated reasoning within mobile device constraints \cite{chen2019deep}. Quantization approaches reduce model memory footprints by representing parameters with lower-precision numerical formats while preserving accuracy. Pruning techniques eliminate redundant network connections to decrease computational requirements. Knowledge distillation transfers capabilities from large teacher models to compact student architectures suitable for edge deployment. These advances collectively enable the on-device processing foundation essential for privacy-preserving personalization.

\textbf{Efficient Fine-Tuning Methods.} Traditional neural network adaptation requires updating all model parameters---computationally prohibitive for on-device execution. Low-Rank Adaptation (LoRA) introduced a breakthrough approach that freezes pretrained weights while training only compact rank decomposition matrices injected at each transformer layer \cite{hu2022lora}. This technique reduces trainable parameters by factors exceeding 10,000$\times$ compared to full fine-tuning while achieving comparable or superior task performance. LoRA's efficiency characteristics align precisely with on-device personalization requirements: minimal memory overhead, rapid adaptation cycles, and the ability to maintain multiple task-specific adaptations simultaneously. Aria extends these principles for continuous, incremental user personalization within mobile device constraints.

\textbf{Privacy-Preserving Distributed Learning.} Federated learning established that machine learning models could achieve competitive performance while keeping training data decentralized across edge devices \cite{mcmahan2017communication}. Subsequent research integrated differential privacy guarantees with federated protocols, providing formal privacy protections against inference attacks on model updates. These developments demonstrate the technical feasibility of learning from sensitive data without centralization. Aria builds upon these foundations while moving beyond the server-aggregation paradigm toward fully device-contained personalization that eliminates even model update transmission.

\textbf{User Modeling and Personalization.} Research on user modeling for information systems and recommender systems has evolved from simple preference profiles toward sophisticated representations capturing behavioral patterns, contextual influences, and temporal dynamics \cite{eke2019survey}. Contemporary approaches employ deep learning architectures to model complex user behavior sequences and multi-type interactions \cite{he2023survey}. The field increasingly recognizes that users exhibit different preferences and behaviors across contexts, motivating research into context-dependent modeling. Aria's Sims Technology synthesizes these advances into dynamic digital representations that evolve through ongoing interaction rather than static profile snapshots.

\textbf{Context-Aware Computing.} Foundational work defined context as any information characterizing the situation of an entity relevant to user-application interaction, establishing that systems could adapt behavior based on location, time, identity, and activity \cite{dey2001understanding}. Subsequent research developed context modeling and reasoning techniques supporting context-aware application development \cite{bettini2010survey}. The proliferation of sensors in mobile devices---accelerometers, GPS, ambient light, proximity---provides rich signals for inferring user context without explicit input. Aria leverages these advances for automatic persona selection and blending based on detected contextual state.

\vspace*{-8pt}
\subsection{Proposed End Goals and Translation Path}
\vspace*{-8pt}
This project pursues specific translation outcomes that collectively establish Aria as a commercially viable product addressing demonstrated market needs.

\textbf{Translation Outcome 1: Functional Prototype.} The primary deliverable is a deployable Aria application demonstrating Sims Technology on iOS and Android platforms. This prototype will implement core capabilities including on-device persona creation and evolution, automatic context detection and persona selection, privacy-preserving personalization without cloud data transmission, and cross-context persona synthesis. The prototype will undergo iterative refinement through beta testing with target user populations, incorporating customer discovery insights from I-Corps training.

\textbf{Translation Outcome 2: Startup Formation.} NxTAGE will incorporate as a commercial entity positioned to pursue market deployment following project completion. Corporate formation activities include establishing appropriate legal structure, formalizing intellectual property assignments, developing governance frameworks, and preparing investor materials for follow-on funding pursuit.

\textbf{Translation Outcome 3: Strategic Partnership Development.} Customer discovery activities will identify and cultivate relationships with potential strategic partners including device manufacturers seeking privacy-differentiated AI capabilities, enterprise software providers requiring compliant AI assistance integration, and platform operators interested in licensing Sims Technology for their ecosystems.

\textbf{Translation Outcome 4: Intellectual Property Protection.} The project will pursue patent protection for core innovations including the Sims Technology architecture, on-device persona synthesis methods, context-aware persona selection algorithms, and privacy-preserving fine-tuning adaptations. A provisional patent application has been filed; this project will support prosecution toward issuance and strategic continuation filings.

\textbf{Translation Outcome 5: Follow-on Funding Position.} Project completion will position NxTAGE for NSF SBIR/STTR Phase I application, venture capital investment, or strategic partnership investment. Demonstrated technical feasibility, validated market demand through customer discovery, and established intellectual property will collectively strengthen subsequent funding applications.

\vspace*{-8pt}
\subsection{Potential Value and Interested Stakeholders}
\vspace*{-8pt}
\textbf{Market Opportunity.} The intelligent virtual assistant market exhibits substantial scale and growth trajectory. Industry analyses project the global market to reach approximately \$14 billion by 2030, representing a compound annual growth rate exceeding 24\% \cite{grandview2023iva}. The market expansion reflects increasing consumer adoption of smart devices, enterprise deployment of conversational AI for customer service and productivity enhancement, and integration of AI assistants across automotive, healthcare, and smart home applications. Critically, privacy-focused segments demonstrate accelerated growth as consumers and enterprises increasingly prioritize data protection, creating specific demand for solutions addressing the privacy-personalization paradox.

\textbf{Target Market Segments.} Aria addresses three primary market segments with distinct value propositions. The \textit{privacy-conscious consumer segment} comprises early adopters who recognize privacy value and will pay premium prices for data sovereignty. This segment provides initial market validation and generates testimonials supporting broader adoption. The \textit{enterprise compliance segment} includes organizations in regulated industries---healthcare, financial services, legal, government---that cannot deploy cloud-dependent AI assistants due to data handling restrictions. Aria enables these organizations to realize AI productivity benefits while maintaining regulatory compliance. The \textit{device manufacturer segment} represents original equipment manufacturers seeking differentiation through privacy-respecting AI capabilities integrated at the operating system or hardware level. Strategic partnerships with device manufacturers offer licensing revenue and scaled distribution channels.

\textbf{Competitive Positioning.} Aria occupies a unique competitive position combining complete on-device processing (unlike cloud-dependent incumbents), multi-context persona modeling (unlike single-profile alternatives), and dynamic evolutionary adaptation (unlike static preference systems). This positioning addresses unmet market needs while establishing defensible differentiation through patented technology and accumulated user trust.

\textbf{Economic Impact.} Successful commercialization generates economic value through direct employment at NxTAGE, technology licensing revenue contributing to university technology transfer outcomes, downstream economic activity from enterprise productivity improvements enabled by compliant AI assistance, and potential acquisition value if strategic acquirers recognize Aria's market position and technology assets.

%%%%%%%%%%%%%%%%%%%%%%%%%%%%%%%%%%%%%%%%%%%%%%%%%%%%%%%%%%%%%%%%%%%%%%%%%%%%%%%
% Section 3: Development of Solutions
% NSF TTP-T Proposal: NxTAGE/Aria Privacy-First Personal AI Assistant
% Target length: 5-6 pages
%%%%%%%%%%%%%%%%%%%%%%%%%%%%%%%%%%%%%%%%%%%%%%%%%%%%%%%%%%%%%%%%%%%%%%%%%%%%%%%

\vspace*{-8pt}
\section{Development of Solutions}
\label{sec:development}
\vspace*{-8pt}
This section presents the technical development plan for translating Aria's privacy-first personal AI assistant from research prototype to market-ready product. We address the fundamental engineering challenges of delivering sophisticated personalization entirely on-device, describe our iterative validation methodology, establish measurable milestones with decision criteria, and outline the path to commercial deployment.

\vspace*{-8pt}
\subsection{Technical Barriers and Anticipated Hurdles}
\vspace*{-8pt}
Realizing on-device personalized AI presents four interconnected technical challenges that require novel solutions combining advances in model compression, context-aware computing, privacy-preserving machine learning, and human-computer interaction.

\textbf{Challenge 1: On-Device Computational Constraints.}
Contemporary mobile devices impose severe constraints on memory, processing power, and energy consumption that fundamentally limit the deployment of sophisticated AI models. State-of-the-art language models require gigabytes of memory and substantial computational resources, creating a fundamental tension with mobile deployment requirements \cite{han2016deep}. 

Our approach employs a multi-faceted compression strategy combining complementary techniques. First, we apply structured pruning to systematically remove redundant network connections while preserving model accuracy. Research demonstrates that neural networks contain significant redundancy, with pruning achieving 9-13$\times$ reduction in connections without accuracy loss \cite{liang2021pruning}. Second, we implement quantization to reduce numerical precision from 32-bit floating point to 8-bit or 4-bit integers, enabling 4-8$\times$ memory reduction with minimal accuracy degradation \cite{gholami2022survey}. Third, we employ knowledge distillation to transfer capabilities from large teacher models to compact student models optimized for mobile deployment \cite{hinton2015distilling}. The distillation process preserves the teacher's generalization capabilities while dramatically reducing computational requirements.

For personalization, we adopt Low-Rank Adaptation (LoRA), which reduces trainable parameters by approximately 10,000$\times$ compared to full fine-tuning while achieving comparable adaptation quality \cite{hu2022lora}. This enables continuous on-device learning without prohibitive memory requirements. We target deployment on devices with MobileNet-class architectures, which achieve strong performance using depthwise separable convolutions that reduce computational cost by 8-9$\times$ compared to standard convolutions \cite{howard2017mobilenets}.

Risk mitigation includes tiered functionality based on device capabilities, with graceful degradation for resource-constrained devices, and optional edge computing integration for computationally intensive operations when privacy-preserving protocols can be maintained.

\textbf{Challenge 2: Multi-Context Persona Synthesis.}
Users inhabit multiple life contexts---professional, personal, family, social---each requiring different interaction patterns and information access. Current AI assistants fail to recognize these context boundaries, leading to inappropriate responses and privacy violations through context collapse \cite{dey2001understanding}. 

Sims Technology addresses this through dynamic digital representations that evolve based on contextual signals. Our approach leverages smartphone sensor data for context inference, building on advances in human activity recognition that achieve over 90\% accuracy in distinguishing activities using accelerometer and gyroscope data \cite{wang2019deep}. We extend these techniques to recognize higher-level contexts (work meetings, family dinner, exercise) through multi-modal sensor fusion combining location patterns, temporal signals, application usage, and ambient audio characteristics.

Context switching employs a hierarchical state machine that manages transitions between personas while maintaining appropriate information boundaries. Each Sim maintains independent preference models, communication styles, and access permissions, enabling the assistant to adapt its behavior appropriately without explicit user intervention. User override mechanisms ensure transparency and control, with clear visual indicators showing the currently active context.

\textbf{Challenge 3: Privacy-Preserving Personalization.}
Achieving deep personalization traditionally requires extensive data collection and cloud-based processing, creating fundamental privacy vulnerabilities. Our approach must deliver equivalent personalization quality while ensuring all sensitive data remains exclusively on-device.

We implement differential privacy guarantees for any model updates, building on foundational work demonstrating that deep neural networks can be trained under modest privacy budgets without unacceptable accuracy degradation \cite{abadi2016deep}. The differentially private stochastic gradient descent (DP-SGD) algorithm adds calibrated noise during training, providing formal mathematical guarantees that individual user data cannot be extracted from model parameters.

For scenarios requiring aggregated insights across users (e.g., improving base model performance), we employ federated learning protocols that enable collaborative model improvement without centralizing raw data \cite{mcmahan2017communication}. Secure aggregation cryptographic protocols ensure that only aggregated gradients---not individual contributions---are visible to any coordinating server \cite{bonawitz2017practical}.

Hardware-level protection utilizes Trusted Execution Environments (TEEs), specifically ARM TrustZone on mobile devices, to isolate sensitive computations from the main operating system \cite{pinto2019demystifying}. TEEs provide hardware-enforced memory isolation ensuring that even a compromised operating system cannot access protected data or computations. We designate TEE enclaves for credential storage, biometric processing, and cryptographic operations underlying Sims security.

\textbf{Challenge 4: User Trust and Adoption.}
Privacy-focused products face a credibility challenge: users have been conditioned by repeated privacy violations to distrust corporate privacy claims. Research on technology acceptance demonstrates that perceived usefulness and ease of use are necessary but insufficient for adoption of AI technologies; trust and transparency emerge as critical additional factors \cite{davis1989perceived, sohn2020technology}.

Our trust-building strategy operates on multiple levels. Technical transparency includes open-sourcing core privacy-preserving components, enabling independent security audits, and publishing formal privacy proofs for our differential privacy implementations. User-facing transparency provides real-time dashboards showing exactly what data is stored locally, how it influences personalization, and explicit confirmation that no data leaves the device. Progressive disclosure gradually introduces AI capabilities as users develop comfort, avoiding the ``uncanny valley'' of unexpectedly sophisticated behavior.

We conduct iterative user studies following established human-computer interaction methodologies to validate that our transparency mechanisms effectively build trust without overwhelming users with technical complexity \cite{hussain2009current}. Findings will inform interface refinements ensuring privacy guarantees are both robust and comprehensible.

\vspace*{-8pt}
\subsection{Testing, Refinement, and Validation Plan}
\vspace*{-8pt}
Our development follows an iterative methodology integrating agile software development with user-centered design principles \cite{silva2011user}. This approach enables rapid iteration on technical components while maintaining continuous alignment with user needs and market requirements.

\textbf{Phase 1: Core Technology Validation (Months 1-12).}
The first year focuses on establishing technical feasibility for on-device personalization at production quality levels.

\textit{Algorithm Development and Benchmarking:} We develop and rigorously benchmark on-device fine-tuning algorithms targeting sub-500ms latency for inference and sub-2GB memory footprint for the complete Sims model. Benchmarking employs standardized datasets and comparison against cloud-based alternatives to quantify any accuracy-efficiency tradeoffs.

\textit{Sims Framework Prototype:} We implement the initial Sims representation framework, including context detection algorithms, persona state management, and preference learning mechanisms. Technical validation confirms that multi-context switching operates reliably with $>$85\% context detection accuracy.

\textit{Controlled User Studies:} We conduct user studies with approximately 50 participants following IRB-approved protocols. Studies employ mixed methods combining quantitative usability metrics (task completion time, error rates, System Usability Scale scores) with qualitative feedback through semi-structured interviews. Participant recruitment emphasizes demographic diversity to identify potential usability issues across user populations.

\textit{I-Corps Customer Discovery:} Concurrent with technical development, the project team completes NSF I-Corps training, conducting 100+ customer discovery interviews to validate market assumptions and refine the value proposition. Findings directly inform Phase 2 development priorities.

\textit{Deliverables:} Technical feasibility demonstration; validated fine-tuning algorithms; preliminary Sims framework; customer discovery insights; I-Corps completion certificate.

\textbf{Phase 2: Prototype Development (Months 13-24).}
Year two translates validated algorithms into functional mobile applications while expanding user validation.

\textit{Cross-Platform Development:} We develop functional prototypes for both iOS and Android platforms, ensuring feature parity and consistent user experience. Development employs native platform capabilities for optimal performance, particularly for sensor access and TEE integration.

\textit{Multi-Context Implementation:} We implement complete multi-context persona switching with the full Sims Technology stack. Integration testing validates seamless transitions between contexts and appropriate information boundary enforcement.

\textit{Expanded User Studies:} We expand user studies to approximately 200 participants across diverse demographics, including targeted recruitment of privacy-conscious users and enterprise information workers representing key market segments. Studies employ longitudinal designs with 4-8 week usage periods to assess sustained engagement and trust development.

\textit{Third-Party Security Audit:} We engage independent security researchers to audit privacy implementations, including penetration testing, code review of TEE integrations, and formal verification of differential privacy bounds.

\textit{Deliverables:} Beta-ready iOS and Android prototypes; security audit report; expanded user study findings; refined product specifications.

\textbf{Phase 3: Market Validation (Months 25-36).}
The final year validates product-market fit through real-world deployment and establishes commercialization foundations.

\textit{Beta Deployment:} We deploy beta versions to approximately 1,000 early adopters recruited through I-Corps networks, privacy advocacy communities, and enterprise pilot partners. Beta participants provide continuous feedback through in-app mechanisms and periodic surveys.

\textit{Usage Analytics:} Privacy-preserving on-device analytics capture engagement metrics (session frequency, feature usage, retention) without compromising user privacy. Aggregated, differentially private summaries inform product refinement.

\textit{Enterprise Pilots:} We conduct structured pilots with 2-3 enterprise partners in privacy-sensitive industries (healthcare, legal, financial services) to validate B2B value proposition and identify enterprise-specific requirements.

\textit{Commercialization Planning:} Based on market validation findings, we finalize go-to-market strategy, pricing models, and partnership structures. Preparation of materials for follow-on funding (SBIR/STTR, venture capital) proceeds in parallel.

\textit{Deliverables:} Market-ready product; validated product-market fit metrics; enterprise pilot results; commercialization plan; investor materials.

\vspace*{-8pt}
\subsection{Development Timeline and Milestones}
\vspace*{-8pt}
Table \ref{tab:milestones} presents the 36-month milestone schedule with quarterly deliverables, success metrics, and go/no-go decision criteria.

\begin{table}[h]
\centering
\small
\begin{tabular}{p{0.6cm}p{4.5cm}p{3.2cm}p{4.2cm}}
\toprule
\textbf{Q} & \textbf{Milestone} & \textbf{Success Metric} & \textbf{Go/No-Go Criteria} \\
\midrule
\multicolumn{4}{l}{\textit{Year 1: Core Technology Validation}} \\
1 & On-device inference baseline & Latency $<$500ms (95th \%ile) & Technical feasibility confirmed \\
2 & Fine-tuning algorithm v1.0 & Memory $<$2GB peak & Mobile deployment viable \\
3 & Sims framework prototype & User study approval (n=50) & Technical approach validated \\
4 & I-Corps training completion & 100+ discovery interviews & Market understanding established \\
\midrule
\multicolumn{4}{l}{\textit{Year 2: Prototype Development}} \\
5 & iOS prototype alpha & 100 test users onboarded & Usability threshold met (SUS$>$68) \\
6 & Android prototype alpha & Cross-platform parity & Platform strategy confirmed \\
7 & Multi-context switching & Context accuracy $>$85\% & Core differentiator functional \\
8 & Security audit complete & No critical vulnerabilities & Trust claims independently verified \\
\midrule
\multicolumn{4}{l}{\textit{Year 3: Market Validation}} \\
9 & Public beta launch & 500+ active weekly users & Initial traction demonstrated \\
10 & Enterprise pilot initiation & 2+ signed pilot agreements & B2B viability signal \\
11 & Investor readiness & Pitch deck; financial model & Funding pathway established \\
12 & Commercialization plan & Complete GTM strategy & Sustainable business model defined \\
\bottomrule
\end{tabular}
\caption{36-month milestone timeline with success metrics and decision criteria}
\label{tab:milestones}
\end{table}

\vspace*{-8pt}
\subsection{Success Metrics and Pivot Criteria}
\vspace*{-8pt}
We establish quantitative metrics across technical performance, user engagement, and market validation dimensions, with explicit thresholds triggering strategic pivots if primary approaches prove unviable.

\textbf{Technical Performance Metrics:}
\begin{itemize}[noitemsep]
    \item \textit{Inference Latency:} Target $<$500ms at 95th percentile for natural interaction flow
    \item \textit{Memory Footprint:} Target $<$2GB for complete Sims model enabling broad device compatibility
    \item \textit{Personalization Quality:} Target $>$80\% user satisfaction with assistant responses (validated via periodic surveys)
    \item \textit{Context Detection Accuracy:} Target $>$85\% correct context identification across standard scenarios
    \item \textit{Battery Impact:} Target $<$5\% additional daily battery consumption during typical usage
\end{itemize}

\textbf{User Engagement Metrics:}
\begin{itemize}[noitemsep]
    \item \textit{Retention:} Target $>$40\% weekly active users among beta cohort
    \item \textit{Session Frequency:} Target $>$3 sessions per user per day indicating habitual usage
    \item \textit{Net Promoter Score:} Target NPS $>$50 indicating strong user advocacy
    \item \textit{Trust Score:} Target $>$4.0/5.0 on validated trust measurement instruments
\end{itemize}

\textbf{Market Validation Metrics:}
\begin{itemize}[noitemsep]
    \item \textit{Enterprise Interest:} Target $>$5 qualified enterprise leads by Month 30
    \item \textit{Pilot Conversion:} Target $>$50\% of pilots converting to paid engagements
    \item \textit{Willingness to Pay:} Validate \$5-10/month consumer pricing through conjoint analysis
\end{itemize}

\textbf{Pivot Triggers and Alternative Strategies:}
Strategic pivots activate if primary approaches fail to meet threshold criteria:

\textit{If on-device performance is unachievable:} Pivot to privacy-preserving edge computing architecture using homomorphic encryption or secure multi-party computation for computationally intensive operations while maintaining local data storage.

\textit{If consumer adoption is slow:} Pivot to enterprise-first strategy focusing on regulated industries (healthcare, legal, financial) where privacy requirements create stronger value proposition and longer sales cycles are acceptable.

\textit{If multi-context Sims complexity is prohibitive:} Simplify to single-context personalization with manual context switching, reducing technical complexity while preserving core privacy differentiation.

\textit{If trust barriers persist despite transparency:} Pursue third-party privacy certification (e.g., TRUSTe, Privacy Shield successors) and open-source core components to enable community verification.

\vspace*{-8pt}
\subsection{Risk Assessment and Mitigation}
\vspace*{-8pt}
Table \ref{tab:risks} presents comprehensive risk analysis across technical, market, operational, and regulatory dimensions.

\begin{table}[h]
\centering
\small
\begin{tabular}{p{2.8cm}p{1.3cm}p{1.3cm}p{6.8cm}}
\toprule
\textbf{Risk Category} & \textbf{Likelihood} & \textbf{Impact} & \textbf{Mitigation Strategy} \\
\midrule
\multicolumn{4}{l}{\textit{Technical Risks}} \\
Hardware limitations prevent on-device deployment & Medium & High & Tiered functionality; edge computing fallback; continuous benchmarking on latest devices \\
Model compression degrades quality unacceptably & Medium & High & Ensemble of compression techniques; task-specific optimization; quality gates at each milestone \\
Context detection accuracy insufficient & Medium & Medium & Multi-modal sensor fusion; user feedback loop for continuous improvement; manual override option \\
\midrule
\multicolumn{4}{l}{\textit{Market Risks}} \\
Major platform (Apple/Google) launches competing feature & High & High & Patent protection; first-mover brand establishment; enterprise focus less vulnerable to consumer platform competition \\
Privacy fatigue limits user willingness to engage & Medium & Medium & Simplified messaging; progressive disclosure; third-party validation \\
Enterprise sales cycles exceed runway & Medium & High & Balanced consumer/enterprise strategy; milestone-based enterprise engagement \\
\midrule
\multicolumn{4}{l}{\textit{Operational Risks}} \\
Key personnel departure & Low & High & Knowledge documentation; team redundancy; competitive compensation \\
Development delays & Medium & Medium & Agile methodology with sprint buffers; prioritized feature backlog; regular milestone reviews \\
\midrule
\multicolumn{4}{l}{\textit{Regulatory Risks}} \\
Privacy regulations increase compliance burden & Low & Medium & Privacy-by-design architecture exceeds current requirements; regulatory monitoring; flexible architecture \\
Platform policy changes restrict functionality & Low & High & Multi-platform strategy; native platform relationships; contingency for web-based delivery \\
\bottomrule
\end{tabular}
\caption{Risk assessment matrix with likelihood, impact, and mitigation strategies}
\label{tab:risks}
\end{table}

\vspace*{-8pt}
\subsection{Integration and Market Deployment Strategy}
\vspace*{-8pt}
Our commercialization strategy follows a staged approach balancing market validation with sustainable growth.

\textbf{Phase 1: Direct-to-Consumer Foundation (Months 30-42).}
Initial market entry targets privacy-conscious consumers through direct mobile app distribution. The freemium model offers core privacy-preserving assistant functionality at no cost, with premium subscription (\$7.99/month) unlocking advanced Sims capabilities including unlimited context personas, enhanced personalization depth, and priority feature access. This approach builds user base and brand recognition while validating willingness to pay.

Customer acquisition emphasizes organic channels aligned with privacy-conscious audiences: privacy-focused media coverage, security researcher endorsements, and community building in privacy advocacy spaces. Paid acquisition supplements organic growth with careful attention to customer acquisition cost relative to lifetime value.

\textbf{Phase 2: Enterprise Licensing (Months 36-48).}
Enterprise expansion targets organizations in regulated industries where privacy requirements create acute pain points with existing solutions. The enterprise offering includes dedicated deployment support, administrative controls, compliance documentation, and service level agreements. Pricing follows per-seat licensing (\$15-25/user/month) with volume discounts for large deployments.

Initial enterprise verticals include healthcare (HIPAA compliance requirements), legal (attorney-client privilege protection), and financial services (regulatory data protection mandates). Enterprise sales leverage I-Corps network connections and pilot relationship conversions.

\textbf{Phase 3: OEM Partnerships (Months 42+).}
Long-term strategy includes licensing Sims Technology to device manufacturers seeking privacy differentiation. OEM partnerships provide recurring licensing revenue and dramatically expanded user reach while device manufacturers gain competitive differentiation in increasingly privacy-conscious markets. Partnership structures may include per-device licensing fees, revenue sharing on premium features, or strategic equity investments.

\textbf{Competitive Positioning:}
Aria's competitive moat comprises multiple reinforcing elements. Patent protection covers core Sims Technology innovations including multi-context persona synthesis, on-device evolutionary adaptation, and privacy-preserving personalization methods. Technical barriers to imitation include the substantial R\&D investment required to replicate on-device performance. Network effects emerge as users develop personalized Sims that become increasingly valuable over time and cannot be transferred to competing platforms. First-mover advantage in privacy-first personal AI establishes brand association that competitors cannot easily replicate.

\vspace*{-8pt}
\subsection{Intellectual Property Landscape}
\vspace*{-8pt}
\textbf{Current IP Status:}
The project has established foundational IP protection through a provisional patent application filed in 2024: ``System and Method for Privacy-Preserving Personalized AI Assistance Using Dynamic Context-Aware Digital Representations.'' The provisional application covers core innovations including the Sims Technology architecture, on-device fine-tuning optimizations, multi-context persona management, and privacy-preserving personalization methods.

Additionally, trade secret protections cover implementation details including specific model architectures, training procedures, and optimization techniques that provide competitive advantage but are better protected through secrecy than patent disclosure.

\textbf{IP Strategy:}
Year 1 priorities include converting the provisional patent to utility patent applications covering core innovations, with claims structured to maximize defensive value while enabling licensing flexibility. We budget up to \$50,000 for patent prosecution costs as permitted under NSF TTP-T guidelines.

Freedom-to-operate analysis will map the existing patent landscape in personal AI assistants, on-device machine learning, and privacy-preserving computation. This analysis identifies potential infringement risks and informs both technical design decisions and partnership/licensing strategies. Key patent holders in adjacent spaces include major technology companies (Apple, Google, Amazon, Microsoft) with substantial voice assistant portfolios, as well as specialized AI companies and academic institutions.

\textbf{Preliminary Landscape Assessment:}
Initial patent landscape review identifies several relevant patent clusters: voice assistant interaction patterns (heavily patented by major platforms), general machine learning optimization techniques (broad prior art limits patentability), and federated learning methods (emerging area with increasing patent activity). Aria's Sims Technology occupies a distinctive position combining on-device personalization, multi-context awareness, and privacy-preserving adaptation in ways not clearly anticipated by existing patents. The freedom-to-operate analysis in Year 1 will provide definitive assessment of infringement risks and design-around requirements.

Defensive publication strategy supplements patent protection by establishing prior art for innovations we choose not to patent, preventing competitors from obtaining blocking patents on complementary techniques.

\vspace*{-8pt}
\section{Project Team}
\label{sec:team}

\vspace*{-8pt}
\subsection{Principal Investigator}
\vspace*{-8pt}
\textbf{Chirag Shah}, Professor, Information School, University of Washington
\begin{itemize}[noitemsep]
    \item Expertise in [relevant areas: information retrieval, HCI, AI/ML]
    \item Prior NSF funding: [relevant awards]
    \item Track record: [publications, patents, prior commercialization]
    \item Role: Project leadership, technical direction, stakeholder engagement
\end{itemize}

\subsection{Postdoc and Graduate Student Researchers}
\textbf{TBD}, Postdoc, University of Washington
\begin{itemize}[noitemsep]
    \item Research focus: [privacy, user studies, etc.]
    \item Role: Privacy architecture, user study design and execution
\end{itemize}

\noindent
\textbf{Harshita Chopra}, PhD Student, Computer Science, University of Washington
\begin{itemize}[noitemsep]
    \item Research focus: [on-device ML, personalization, etc.]
    \item Role: Core algorithm development, prototype implementation
    \item Entrepreneurial Lead for I-Corps training
\end{itemize}

\subsection{Advisory Support}
\textbf{Technology Commercialization Expert} (required for TTP)
\begin{itemize}[noitemsep]
    \item Name/TBD, Title/Affiliation
    \item Experience: [relevant commercialization experience]
    \item Role: Market strategy guidance, industry mentor for I-Corps
\end{itemize}

% Optional for TTP-T
\subsection{Collaborators/Partners}
\begin{itemize}[noitemsep]
    \item UW CoMotion (technology transfer support)
    \item Industry advisor if applicable
    \item Other collaborators
\end{itemize}

\vspace*{-8pt}
\section{Broader Impacts}
\label{sec:impacts}
\vspace*{-8pt}
This section articulates the broader impacts of the proposed research beyond its 
technical contributions, addressing economic benefits, educational outcomes, 
diversity initiatives, and the post-award commercialization pathway.

\vspace*{-8pt}
\subsection{Economic and Societal Benefits}
\vspace*{-8pt}
\textbf{Economic Impact.}
The proposed work will generate significant economic benefits at regional and 
national levels. Startup formation is expected to create 5--10 high-skilled 
positions within three years post-award, with growth potential to 25--50 
employees within five years of commercial launch. These positions will span 
machine learning engineering, privacy systems architecture, user experience 
design, and business development, contributing to the regional technology 
workforce. The project strengthens Seattle's position as a hub for AI innovation 
and supports regional economic development initiatives. Beyond direct employment, 
the technology has potential to shift the personal AI market toward 
privacy-respecting alternatives, challenging the prevailing surveillance-based 
business models that currently dominate the intelligent virtual assistant space.

\textbf{Societal Impact.}
The proposed Aria system addresses fundamental concerns about individual agency 
in the digital age. Research has demonstrated that data sovereignty encompasses 
meaningful control, ownership, and claims to data articulated by agents ranging 
from individuals to entire societies~\cite{hummel2021data}. By processing all 
personal information on-device, Aria empowers users to maintain genuine control 
over their digital footprint without sacrificing the benefits of AI-powered 
assistance. This approach directly addresses the tension between technological 
convenience and privacy preservation that affects millions of potential users.

Furthermore, the project contributes to rebuilding public trust in AI systems. 
Large-scale surveys indicate that a majority of individuals express wariness 
about trusting AI systems, with privacy and data governance consistently ranked 
among the most important factors for establishing trust~\cite{gillespie2023trust}. 
By demonstrating a viable alternative to cloud-dependent AI assistants, this 
research provides evidence that privacy-preserving personalization is technically 
achievable and commercially viable.

The technology also enables AI assistance in sensitive domains---including 
healthcare, legal services, and financial planning---where privacy concerns 
currently limit adoption. Users in these contexts require assurance that their 
interactions remain confidential, a guarantee that on-device processing can 
provide but cloud-based alternatives cannot.

\vspace*{-8pt}
\subsection{Educational and Training Outcomes}
\vspace*{-8pt}
\textbf{Student Development.}
Graduate students participating in this project will receive comprehensive 
training that bridges technical research and entrepreneurial practice. The 
mandatory I-Corps program provides experiential learning in customer discovery, 
market validation, and business model development---competencies increasingly 
essential for researchers seeking to translate their work into societal 
impact~\cite{duval2020training}. The I-Corps curriculum, which has facilitated 
the launch of over 1,000 startups nationally, uses the Lean LaunchPad methodology 
to help researchers evaluate commercial opportunities through direct engagement 
with potential customers and industry stakeholders.

Beyond formal I-Corps training, students gain hands-on experience in technology 
commercialization, intellectual property strategy, and startup formation. 
Mentorship from industry professionals and technology transfer specialists 
provides exposure to the practical challenges of bringing academic research to 
market. Students who complete the project will be well-positioned to transition 
into startup employment, pursue entrepreneurial ventures of their own, or bring 
commercialization expertise to academic or industry research positions.

\textbf{Broader Educational Contributions.}
Insights from this project will inform course materials in privacy-preserving 
machine learning, mobile systems design, and technology entrepreneurship. 
Undergraduate research opportunities will be created through directed study 
projects and summer research positions, providing early exposure to cutting-edge 
research at the intersection of AI and privacy. Open-source components developed 
during the project---including on-device inference frameworks and privacy 
evaluation tools---will be made available for educational use, enabling other 
institutions to incorporate these materials into their curricula.

\vspace*{-8pt}
\subsection{Broadening Participation}
\vspace*{-8pt}
The project maintains a strong commitment to diversity, equity, and inclusion 
throughout all activities. Research demonstrates that women and racial minorities 
remain significantly underrepresented in STEM fields, particularly in computer 
science and engineering, where systemic barriers continue to limit 
participation~\cite{rincon2023inclusion}. To address these disparities, the 
project implements several targeted initiatives.

Recruitment efforts will actively seek team members from underrepresented groups, 
leveraging partnerships with campus diversity programs and professional 
organizations focused on broadening participation in computing. User studies will 
be explicitly designed to include participants from diverse demographic 
backgrounds, ensuring that the technology meets the needs of varied user 
populations rather than reflecting the preferences of a narrow demographic 
segment. Accessibility considerations will be integrated into product design from 
the earliest stages, following universal design principles that benefit users 
with disabilities while improving usability for all users.

Outreach activities will engage minority-serving institutions for potential 
collaboration, creating pathways for students from these institutions to 
participate in project activities. These efforts align with NSF's broader goals 
of developing a diverse, globally competitive STEM workforce.

\vspace*{-8pt}
\subsection{Responsible AI Development}
\vspace*{-8pt}
The project embraces principles of responsible AI development, recognizing that 
ethical considerations must be integrated throughout the technology lifecycle. 
Researchers have proposed unified frameworks emphasizing beneficence, 
non-maleficence, autonomy, justice, and explicability as foundational principles 
for AI systems that serve societal good~\cite{floridi2019unified}. The Aria 
system operationalizes these principles through its privacy-first architecture, 
which respects user autonomy by keeping personal data under individual control 
and promotes justice by making personalized AI assistance accessible without 
requiring users to surrender their privacy.

Transparency mechanisms will be built into the system to help users understand 
how their Sims Technology profiles influence AI recommendations. This 
explicability supports informed consent and enables users to modify or constrain 
the system's behavior according to their preferences. Regular ethical reviews 
will assess the technology's alignment with responsible AI principles as 
development progresses.

\vspace*{-8pt}
\subsection{Post-Award Commercialization Pathway}
\vspace*{-8pt}
\textbf{Anticipated Next Steps.}
The project establishes a clear pathway from research to commercial impact:

\begin{enumerate}[noitemsep]
    \item \textbf{Startup Formation (Target: Month 24):} Incorporate NxTAGE as a 
    commercial entity, with initial focus on direct-to-consumer market validation.
    
    \item \textbf{Seed Investment (Months 24--30):} Pursue angel and venture 
    capital investment based on demonstrated user traction and validated 
    product-market fit from beta deployments.
    
    \item \textbf{NSF SBIR/STTR (Months 30--36):} Submit Phase I application for 
    continued development funding, leveraging I-Corps completion and 
    commercialization progress as evidence of commercial potential.
    
    \item \textbf{Strategic Partnerships (Months 30--36):} Formalize relationships 
    with device manufacturers, enterprise customers, or complementary technology 
    providers identified through market validation activities.
    
    \item \textbf{Accelerator Programs:} Apply to relevant technology accelerators 
    (e.g., Y Combinator, Techstars) to accelerate growth and expand investor 
    networks.
\end{enumerate}

Success in these activities will position NxTAGE to pursue Series A funding and scale operations to meet market demand for privacy-respecting AI assistance. The combination of technical differentiation, validated market opportunity, and experienced team provides a foundation for sustainable commercial growth while maintaining the privacy-first principles that define the technology's value proposition.

\vspace*{-8pt}
\section{Results of Prior NSF Support}
\label{sec:prior}
\vspace*{-8pt}
{\bf Shah} was supported under \textbf{III-1717488}  (09/17-08/22; \$500k) as PI. The project developed proactive recommendation systems and CIS methodologies directly applicable to SYNAPSE. \textbf{IM:} This research established theoretical foundations for proactive agent behavior \cite{shah2018information} and developed techniques for collaborative task identification \cite{mitsui2018much, mitsui2018broad} that inform the CIN component of SYNAPSE. \textbf{BI:} The project enhanced the Coagmento collaborative search platform \cite{mitsui2018coagmento}, providing infrastructure foundations for the multi-processor evaluation methodologies proposed in SYNAPSE. Additionally, in {\bf III-1910154} (10/19-04/23; \$500k) as Co-PI, Shah developed explanation generation frameworks and studied their effectiveness in human-agent interactions \cite{fu2020fairness, verma2024counterfactual}, contributing essential groundwork for the trust mechanisms and transparency requirements in collaborative intelligence systems. Shah serves as co-PI on \textbf{III-2336769} and \textbf{CPS-2438935} planning grant, both recently initiated projects that complement the proposed project's objectives.

\textbf{Relation to Proposed Work:}
The proposed Aria project directly extends [specific findings] from this 
prior award, translating the fundamental research on [topic] into a 
commercial product addressing [market need].

%%%%%%%%%%%%%%%%%%%%%%%%%%%%%%%%%%%%%%%%%%%%%%%%%%%%%%%%%%%%%%%%%%%%%%%%%%%%%%%
% END OF PROJECT DESCRIPTION (15 pages)
%%%%%%%%%%%%%%%%%%%%%%%%%%%%%%%%%%%%%%%%%%%%%%%%%%%%%%%%%%%%%%%%%%%%%%%%%%%%%%%

%%%%%%%%%%%%%%%%%%%%%%%%%%%%%%%%%%%%%%%%%%%%%%%%%%%%%%%%%%%%%%%%%%%%%%%%%%%%%%%
% SUPPLEMENTARY MATERIALS CHECKLIST
% (These are separate submission components, not part of 15-page limit)
%%%%%%%%%%%%%%%%%%%%%%%%%%%%%%%%%%%%%%%%%%%%%%%%%%%%%%%%%%%%%%%%%%%%%%%%%%%%%%%

% \section*{Submission Checklist}

% \subsection*{Required Components}
% \begin{itemize}
%     \item[$\square$] Cover Sheet (title must start with ``NSF TTP-T:'')
%     \item[$\square$] Project Summary (1 page)
%     \item[$\square$] Project Description (15 pages max)
%     \item[$\square$] References Cited
%     \item[$\square$] Biographical Sketch(es)
%     \item[$\square$] Budget and Budget Justification
%     \item[$\square$] Current and Pending Support
%     \item[$\square$] Facilities, Equipment, and Other Resources
%     \item[$\square$] Data Management and Sharing Plan (2 pages max)
% \end{itemize}

% \subsection*{Optional Components}
% \begin{itemize}
%     \item[$\square$] Letters of Collaboration (if partners)
%     \item[$\square$] Resubmission Change Description (if applicable, 2 pages max)
% \end{itemize}

% \subsection*{Budget Reminders}
% \begin{itemize}
%     \item[$\square$] I-Corps training: \$45,000 on Line G.6, \$5,000 indirect (Year 1)
%     \item[$\square$] PI Meeting travel: \$2,500/person for PI + 1 student (Year 2)
%     \item[$\square$] Patent expenses: Up to \$50,000 total (optional)
%     \item[$\square$] Total budget $\leq$ \$1,200,000
% \end{itemize}

% %%%%%%%%%%%%%%%%%%%%%%%%%%%%%%%%%%%%%%%%%%%%%%%%%%%%%%%%%%%%%%%%%%%%%%%%%%%%%%%
% % BUDGET TEMPLATE
% %%%%%%%%%%%%%%%%%%%%%%%%%%%%%%%%%%%%%%%%%%%%%%%%%%%%%%%%%%%%%%%%%%%%%%%%%%%%%%%


\newpage
\setcounter{page}{1}
\bibliographystyle{plain}
\bibliography{references}
\end{document}